\chapter{Advanced Cryptographic Primitives}
\label{ch:advanced-primitives}

\abstract*{Lattice-based cryptography enables capabilities far beyond encryption and signatures. This chapter introduces fully homomorphic encryption, zero-knowledge proofs, and secure multiparty computation---primitives that were theoretical curiosities until lattice theory made them practical. For physicists, these protocols offer striking parallels: noise accumulation mirrors thermodynamics, zero-knowledge echoes quantum measurement, and multiparty computation resembles distributed quantum systems.}

% =============================================================
% SOURCE: notes.tex §6.2 (lines 8327–8905)
% REUSE: ~40%. Good overviews with physics analogies.
% Keep focused — this is a gateway chapter, not a treatise.
% =============================================================

\section{Fully Homomorphic Encryption}
\label{sec:fhe}

% TODO: Adapt from notes.tex §6.2.1
% - The concept: computing on encrypted data
% - Gentry's blueprint: somewhat HE + bootstrapping
% - LWE-based FHE: BGV, BFV, CKKS, TFHE
% - The noise growth challenge (physics: thermal noise analogy)
% - Bootstrapping: the noise reset operation
% - Practical FHE today: Microsoft SEAL, TFHE-rs, Zama
% - Cite: Gentry 2009, BGV 2012, CKKS 2017

\section{Zero-Knowledge Proofs}
\label{sec:zkp}

% TODO: Adapt from notes.tex §6.2.2
% - The concept: proving without revealing
% - Interactive proofs and the Fiat-Shamir heuristic (callback to Ch. 8)
% - Lattice-based ZK proofs: Stern-type protocols
% - Post-quantum SNARKs and STARKs
% - Applications: privacy-preserving authentication, anonymous credentials
% - Cite: Goldwasser et al. 1989, Stern 1993, Ben-Sasson et al. (STARKs)

\section{Secure Multiparty Computation}
\label{sec:mpc}

% TODO: Adapt from notes.tex §6.2.3
% - The millionaires' problem and general MPC
% - Secret sharing: Shamir's scheme
% - Garbled circuits: Yao's protocol
% - MPC from lattice-based OT (oblivious transfer)
% - Applications: private set intersection, federated learning

\section{Lattice-Based Advanced Constructions}
\label{sec:lattice-advanced}

% TODO: Adapt from notes.tex §6.2.4
% - Attribute-based encryption
% - Identity-based encryption from lattices
% - Threshold signatures and distributed key generation
% - Functional encryption: computing specific functions on encrypted data

%% Exercises
\begin{prob}
\label{prob:ch14-fhe}
In an LWE-based FHE scheme, a fresh ciphertext has noise magnitude $\sigma$. After one homomorphic multiplication, the noise grows to approximately $\sigma^2 \cdot n$ where $n$ is the lattice dimension. If the modulus is $q$ and decryption fails when noise exceeds $q/2$, compute the maximum multiplicative depth as a function of $\sigma$, $n$, and $q$. For $\sigma = 3.2$, $n = 1024$, $q = 2^{27}$, how many multiplications can be performed before bootstrapping is required?
\end{prob}
